\begin{resume}

% \newcommand\addengcontentsnonumber[2]{%
%     \addcontentsline{toe}{#1}{#2}}
% \addengcontentsnonumber{chapter}{Published papers and research results while pursuing the degree}


%   \researchitem{发表的学术论文} % 发表的和录用的合在一起

%   % 1. 已经刊载的学术论文（本人是第一作者，或者导师为第一作者本人是第二作者）
%   \begin{publications}
%     \item [1]刘英杰，黄嘉琦，姜玉凤，邵宇琪，杨文韬，于紫凝，郑海永.融合电磁和地声特征的地震预测集成学习方法[J].计算机技术与发展,2024,34(08):166-174.
%     % \item [2]Guo Z, Liu Y, Zhang J, et al. Face Forgery Detection with Elaborate Backbone[J]. arXiv preprint arXiv:2409.16945, 2024.
%   \end{publications}
% % \begin{publications}
% %     \item XXX. Intrinsic Image Harmonization[C]//Proceedings of the IEEE/CVF Conference on Computer Vision and Pattern Recognition. 2021: 16367-16376.（EI检索，CCF-A类会议）（第一作者）
% %     \item XXX. Image Harmonization with Transformer[C]//Proceedings of the IEEE/CVF International Conference on Computer Vision. 2021.（已录用，EI检索，CCF-A类会议）（第一作者）
% %     \item XXX. Few-shot Fish Image Generation and Classification[C]//Global Oceans 2020: Singapore–US Gulf Coast. IEEE, 2020: 1-6.（海洋科学技术A类会议）（第一作者）
% %     \item XXX. Underwater Image Enhancement Based on Intrinsic Images[C]//Global Oceans 2021.（已录用，海洋科学技术A类会议）（第一作者）
% %   \end{publications}

%   % 2. 尚未刊载，但已经接到正式录用函的学术论文（本人为第一作者，或者
%   %    导师为第一作者本人是第二作者）。
% %   \begin{publications}[before=\publicationskip,after=\publicationskip]
% %     % \item XXX X, XXX X, XXX X, et al., Title XXX, IEEE XXX, vol. X, pp. XXX-XXX, 20XX.（SCI/EI XXX）
% %   \end{publications}

%   % 3. 其他学术论文。可列出除上述两种情况以外的其他学术论文，但必须是
%   %    已经刊载或者收到正式录用函的论文。
% %   \begin{publications}
% %     % \item XXX X, XXX X, XXX X, et al., Title XXX, IEEE XXX, vol. X, pp. XXX-XXX, 20XX.（SCI/EI XXX）
% %   \end{publications}

%   % \researchitem{研究成果} % 有就写，没有就删除
%   % \begin{achievements}
%   %   \item 刘英杰, 黄嘉琦等. XXX挑战赛X等奖, 20XX年. 主办方：XXX，XXX，XXX.
%   %   \item 刘英杰, 黄嘉琦等. XXX挑战赛X等奖, 20XX年. 主办方：XXX，XXX，XXX.
%   %   \item 刘英杰, 黄嘉琦等. XXX挑战赛X等奖, 20XX年. 主办方：XXX，XXX，XXX.
%   % \end{achievements}
% ==========================================================
% \vspace{32bp} % 段前24磅
% \begin{center}  % 居中无缩进
%     \sanhao\heiti 攻读硕士学位期间取得的研究成果 % 三号，黑体
% \end{center}
% \vspace{14bp} % 段后18磅

% ==========================================================
\newcommand{\smalltitle}[1]{%
    \noindent % 无缩进
    {\heiti\zihao{4} #1} % 黑体四号字，左对齐
}
\addengcontentsnonumber{chapter}{Published papers and research results while pursuing the degree}
% ------------------------------------------------------
% \newcommand{\researchitem}[1]{%
%   \vspace{32bp}{\sihao\heiti\centerline{#1}}\par\vspace{14bp}}

% ------------------------------------------------------
\smalltitle{一、学术论文}
\begin{achievements}
  \item Face Forgery Detection via Unraveling Temporal Forgery Cues. IEEE/CVF Conference on Computer Vision and Pattern Recognition (CVPR), 2025. 共同一作(唯一学生作者).（已接收，CCF-A类会议，对应本文第四章工作）

\end{achievements}

\smalltitle{二、科研项目} 
% ------------------------------------------------------
\begin{achievements}[start=1]
        \item 中央高校基本科研业务费专项，研究生自主科研项目，人脸面部属性和外观特征分离协同的视觉鉴伪，202461014，2023-12-31 至 2024-12-31，1万元，已结题，主持

\end{achievements}

% \begin{achievements}
%         \item 项目名。 2023年度中国海洋大学研究生自主科研项目（\textbf{主持}）。获批并主持该研究生自主科研项目，本年度信息学部仅3名硕士获批。该科研项目旨在通过分离人脸外观与属性特征，并联合外观和属性伪造特征进行鉴伪。
%         \item 成像模型与表示学习联合的伪造人脸检测。国家自然科学基金面上项目（\textbf{课题骨干}）。作为课题骨干参与国家自然科学基金面上项目。该科研项目旨在通过成像模型与表示学习联合，提高模型的人脸表征能力以更好地挖掘人脸伪造特征，提升模型鉴伪的泛化性。
% \end{achievements}

% ------------------------------------------------------
\smalltitle{三、获奖情况}
% ------------------------------------------------------
\begin{achievements}[start=1]
        \item 2025年度研究生国家奖学金
     
\end{achievements}
% ------------------------------------------------------
\end{resume}



